\subsection{Systolic Array Generator}

%Data reuse occurs within the PE array, avoiding broadcasting and resulting in low fan-ins. Systolic array achieves state-of-the-art performance in many domains, like machine learning~\cite{TPU}, due to its high frequency. \autoref{fig:systolic}(a) depicts an example of a systolic array accelerator for matrix multiplication, in which two input matrices are reused through the interconnection between PE array. Generating hardware written in RTL programs is difficult due to the low level of abstraction. However, generating a systolic array at a high-level abstraction is still complicated. Because the HLS procedure often ignores the architecture information, inserting suitable directives is required for an HLS implementation, which is tough even for an expert~\cite{autoSA}. 

Systolic array is a class of accelerators where inputs or results pass through the processor elements (PEs)~\cite{polySA, autoSA, tensorlib, ems}. Despite the fact that systolic array offers good performance, it's challenging to conveniently generate a customized accelerator. In this case study, we use Hector to build a systolic array generator. This generator can generate PE arrays of arbitrary size given a PE's configuration, which can be written in either \tor or \hec. In addition, \texttt{SCF} Dialect is also supported by reusing the previous HLS process (\autoref{sec:hls}).

\begin{figure}[htbp!]
  \centering
  \includegraphics[width=\linewidth]{figure/code.pdf}
  \caption{A 2$\times$n by n$\times$2 matrix multiplication systolic array described in \hec. The interconnection between PEs is described as \texttt{hec.assign}.}
  \label{fig:systolic}
  \vspace{-9pt}
\end{figure}

\textbf{Implementation.} We use an elastic implementation that allows each PE to stall if the data does not arrive or the subsequent PE is not ready to receive new data. Supporting dynamic behavior is critical, particularly when bandwidth is limited. When sufficient bandwidth is ensured, an elastic systolic array performs similarly to the static systolic array with the exception of a few more signals. This dynamic behavior is implemented as an elastic interconnection between PE array. \autoref{fig:systolic} shows a 2$\times$2 systolic array represented in \hec. The outer module is a dynamic style component that enables the implementation to be elastic. All PEs are instantiated from the corresponding component, and data movement between PEs is explicitly stated as \verb|hec.assign|.

A single PE can be configured at multiple levels. For the PE with simple behavior, \texttt{SCF} Dialect can be automatically lowered using the approach described in \autoref{sec:hls}. \hec can be used in applications with no resource limitations, and \tor is capable of handling complex designs. This once again demonstrates the advantages of multi-level IR. Pipelined PE can be easily expressed at these levels due to the explicit pipeline in both levels.

\begin{figure}[htbp!]
  \centering
  \includegraphics[width=\linewidth]{figure/data_systolic.pdf}
  \caption{Cycle count comparison between different implementations of matrix multiplication.}
  \label{fig:data_systolic}
  \vspace{-9pt}
\end{figure}


\textbf{Comparison against Calyx.} Calyx~\cite{} also provides a static systolic array generator for demonstration purposes, but it lacks the expressivity required for hardware generation. More than 3000 lines of code are required to describe an 8×8 systolic array. The fundamental reason is that the control logic is described in a flattened representation that explicitly represents all data moves at each cycle. Another issue in Calyx is that it lacks a pipeline primitive, prohibiting all PEs in generated hardware from being pipelined. In real-world applications, interleaving multiple calculations is commonly used to improve resource utilization ratio, but the lack of pipeline primitives makes this optimization impossible. 

\textbf{Evaluation.} We evaluate matrix-multiplication kernels ranging from 4×4 to 32×32. The PE configuration of the systolic array is written in \texttt{SCF}, which is automatically lowered to \hec. Aside from the systolic arrays generated by Hector and Calyx, we implement a basic program in Vitis HLS that uses pipeline pragma in the inner-most loop. We optimize the HLS implementation by unrolling the outer two loops and partitioning all the matrices in the proper dimension to make a fair comparison. We also evaluate Tensorlib~\cite{tensorlib}, a state-of-the-art systolic array generator, with the same dataflow configuration. The experiment targets Zynq UltraScale+ XCZU3EG FPGA at a 7ns clock period, which is the same as Calyx's setup. We compare the cycle counts of the designs (\autoref{fig:data_systolic}). Hector supports pipeline PE, which further optimizes the performance by interleaving multiple computations. On average, the systolic array generated by Hector improves the performance by 5.6× improvement compared to optimized HLS implementation. In comparison to the HLS implementation, Calyx implementation with sequential PEs can only achieve 79\% performance, demonstrating the importance of pipeline semantics. Hector can achieve the same performance with Tensorlib due to structural description and pipeline support.

The experiment results demonstrate Hector's versatility in applying domain-specific approaches, which is difficult to express in HLS tools. With architectural information, which is allowed in Hector, it's straightforward to outperform a general-purpose HLS tool. The multi-level paradigm also provides the ability to design hardware at various levels.
